\section*{Chapter \ref{ch:g10:floating-point} --- Floating-Point Numbers}

\begin{solution}{exo:g10:floating-point:1}
\code{0.30000000000000004}, \code{False}, \code{True}, \code{True}. The numbers $0.5$, $1.5$, $0.25$ and $0.75$ have a
power of two as denominator ($\tfrac12$, $\tfrac32$, $\tfrac14$, $\tfrac34$): they are stored exactly, and so are the
exact results of these operations, which are floats too. $0.1$ and $0.3$ are not stored exactly.
\end{solution}

\begin{solution}{exo:g10:floating-point:2}
$0.375$: $0.75 \to 0$, $1.5 \to 1$, $1.0 \to 1$, so $0.011$. $0.8125$: $1.625 \to 1$, $1.25 \to 1$, $0.5 \to 0$,
$1.0 \to 1$, so $0.1101$. $5.75 = 101.11$ (the whole part $5$ is $101$, the fraction $0.75$ is $0.11$).
\end{solution}

\begin{solution}{exo:g10:floating-point:3}
$\tfrac13 \to \tfrac23$, digit $0$; $\tfrac23 \to \tfrac43$, digit $1$, keep $\tfrac13$; and $\tfrac13$ has come
back. So $\tfrac13 = 0.010101\dots$ in base $2$, with period $01$, two digits, as the listing printed.
\end{solution}

\begin{solution}{exo:g10:floating-point:4}
\code{True}: $10^{-10} \le 10^{-9} \times 1$. \code{False}: $|a - b| = 10^{-10}$ but $10^{-9} \times 2 \times 10^{-10}$
is far smaller; relatively, $10^{-10}$ and $2 \times 10^{-10}$ differ by a factor of two. \code{True}: the absolute
tolerance $10^{-9}$ is larger than $10^{-10}$.
\end{solution}

\begin{solution}{exo:g10:floating-point:5}
After ten additions \code{x} is \code{0.9999999999999999}, not \code{1.0}, so \code{x != 1.0} stays true and the
loop goes on for ever. Count whole turns instead: \code{for k in range(1, 11): x = k * 0.1}, or test
\code{while x < 0.95}, a bound far from any \omterm{def:g10:floating-point:rounding-error}{rounding error}.
\end{solution}

\begin{solution}{exo:g10:floating-point:6}
The \omterm{def:g10:floating-point:float}{sign} is $-$; the \omterm{def:g10:floating-point:float}{exponent} field is $1024 + 2 = 1026$, so $e = 1026 - 1023 = 3$; the \omterm{def:g10:floating-point:float}{significand} is
$1.11$ in base $2$, that is $1.75$. The number is $-1.75 \times 2^3 = -14$.
\end{solution}

\begin{solution}{exo:g10:floating-point:7}
In \code{exercises/exact_binary.py}: \code{q = Fraction(x)}, then, while \code{q != 0}, double \code{q}, write the
digit \code{1} and subtract $1$ if \code{q >= 1}, else write \code{0}. The loop ends because each doubling removes a
factor $2$ from the denominator. The stored \code{0.1} is $3602879701896397/2^{55}$: its binary digits stop at the
$55$th place.
\end{solution}

\begin{solution}{exo:g10:floating-point:8}
$2^{20}$ lies in $[2^{20}, 2^{21})$, where the gap is $2^{20-52} = 2^{-32} \approx 2.33 \times 10^{-10}$, which is
what \code{math.ulp(2.0**20)} prints. There are $2^{52}$ floats in $[2^{20}, 2^{21})$, as in every interval between
two consecutive powers of two.
\end{solution}

\begin{solution}{exo:g10:floating-point:9}
The first line prints \code{1e+16}: each \code{1.0} is lost in turn, since $10^{16} + 1$ is not a float and rounds back to $10^{16}$.
\code{1.0000000000000002e+16}: $1 + 1 = 2$ is computed first, and $10^{16} + 2$ is a float. \code{sum} also prints
\code{1.0000000000000002e+16}: it keeps the lost digits along the way, like \code{math.fsum}. Adding the small
numbers together first keeps the most digits.
\end{solution}

\begin{solution}{exo:g10:floating-point:10}
\code{k * 0.1} matches \code{k / 10} for $19$ of the $30$ firings, \code{t += 0.1} for only $8$
(\code{exercises/accumulate.py}). The product makes one \omterm{def:g10:floating-point:rounding-error}{rounding error}; the sum carries every error of the
additions before it. That is the second item of the method on keeping \omterm{def:g10:floating-point:rounding-error}{rounding errors} small.
\end{solution}

\begin{solution}{exo:g10:floating-point:11}
\code{return abs(a - b) <= max(rel * max(abs(a), abs(b)), abs_)}, with \code{rel=1e-9} and \code{abs_=0.0} by default
(\code{exercises/almost_equal.py}); the chapter's tests compare it with \code{math.isclose} on $500$ random pairs.
\end{solution}

\begin{solution}{exo:g10:floating-point:12}
$2^{53} + 1$ lies in $[2^{53}, 2^{54})$, where the floats are $m \times 2^{53}$ with $m$ having $52$ binary digits after
the point: they are $2^{53} + 2k$ for whole numbers $k$, all even. $2^{53} + 1$ is odd, so it is not one of them;
$2^{53} + 2$ is ($k = 1$). \code{2**53 + 1.0} gives \code{9007199254740992.0}: the exact sum is exactly halfway between
two floats, and the rule rounds to the one whose last bit is $0$, here $2^{53}$.
\end{solution}

\begin{solution}{exo:g10:floating-point:13}
$16$ values in $[1, 2)$ (\omterm{def:g10:floating-point:float}{exponent} field $3$), $16$ in $[2, 4)$, and $8 \times 16 = 128$ positive values in all, as the
listing printed. Every whole number from $1$ to $31$ is a value: from $1$ to $15$ the gaps are at most $\tfrac12$,
and from $16$ to $31$ the gap is exactly $1$. The smallest whole number that is not a value is $32$, above the
largest value, $31$.
\end{solution}

\begin{solution}{exo:g10:floating-point:14}
With the $+$ \omterm{def:g10:floating-point:float}{sign}, $\sqrt{b^2 - 4} $ is extremely close to $b = 10^8$, and the subtraction $-b + \sqrt{b^2 - 4}$ loses
nearly every digit: Python gives $-7.450580596923828 \times 10^{-9}$. The large root is about $-10^8$, computed with
an addition of two negative numbers and no cancellation; $c$ divided by it gives $-10^{-8}$, which is right to all
printed digits, since the small root is $-10^{-8} - 10^{-24} - \dots$ (\code{exercises/quadratic.py}). The second
method is right; the first is wrong by $25\,\%$.
\end{solution}

\begin{solution}{exo:g10:floating-point:15}
For \code{[nan, 3.0, 7.0]} it returns \code{nan}: \code{best} starts as \code{nan}, and \code{r > nan} is always false.
For \code{[3.0, nan, 7.0]} it returns \code{7.0}, since \code{nan > 3.0} is false and \code{nan} is skipped. The fix
removes the missing readings first, \code{values = [r for r in readings if not math.isnan(r)]}, asserts that some
remain, and takes the largest of them (\code{exercises/nan_max.py}).
\end{solution}

\begin{solution}{pb:g10:floating-point:1}
\textbf{1.} $0.000110011001100110011001100$ (the digits $0011$ repeat from the second place on).

\textbf{2.} Multiply $x$ by $2^{p}$, drop the fractional part, and divide by $2^{p}$ again: what remains is $x$ with
only its first $p$ binary places.

\textbf{3.} The digits from the second place on repeat $0011$; the $24$th and $25$th digits are $1$ and $1$, then
$0011$ again. So the dropped part is $0.11001100\dots \times 2^{-23}$, and $0.1100\,1100\dots$ in base $2$ is
$2 \times 0.0110\,0110\dots = 2 \times 0.4 = 0.8$ (since $0.4$ is $0.1$ doubled twice, $0.0110\,0110\dots$).

\textbf{4.} $0.8 \times 2^{-23} \approx 9.54 \times 10^{-8}$~s per tick, as the program prints.

\textbf{5.} $3600 \times 10 = 36{,}000$ ticks.

\textbf{6.} $36{,}000 \times 9.537 \times 10^{-8} \approx 0.0034$~s, the report's value.

\textbf{7.} $8 \times 36{,}000 = 288{,}000$ ticks, so $288{,}000 \times 9.537 \times 10^{-8} \approx 0.0275$~s, again the
report's value.

\textbf{8.} Each tick adds the same error, so the total error is the number of ticks times the error per tick, and the
number of ticks grows in proportion to the time.

\textbf{9.} $100 \times 36{,}000 = 3{,}600{,}000$ ticks.

\textbf{10.} $3{,}600{,}000 \times 0.8 \times 2^{-23} \approx 0.3433$~s.

\textbf{11.} The program prints $0.0034$, $0.0275$, $0.0687$, $0.1648$, $0.2472$ and $0.3433$~s: the six values of the
report's table, to the four printed decimals (the chapter's tests check them).

\textbf{12.} $0.1287$~s. Only the tenth chopped after $23$ places matches the report: with a $24$-bit register, one
bit must have held something else, so $23$ binary places were left for the fraction.

\textbf{13.} The count restarts at each reboot, so the drift never exceeds its value after $8$ hours, $0.0275$~s.

\textbf{14.} $687 / 0.3433 \approx 2001$, so about $2000$~m/s.

\textbf{15.} $2000 \times 0.0275 = 55$~m, the report's shift.

\textbf{16.} $2000 \times 0.1287 \approx 257$~m: smaller, but still far more than the size of a target.

\textbf{17.} Keep the count of ticks as a whole number, which is exact, and convert it to seconds only when needed with
one multiplication or, better, one division by $10$. The error is then a single rounding, whatever the running time:
items 1 and 2 of the method.

\textbf{18.} The \omterm{def:g10:floating-point:rounding-error}{rounding error} of \code{0.1} contributes $3{,}600{,}000 \times 5.6 \times 10^{-18} \approx 2 \times
10^{-11}$~s, and the rounding of the product, near $360{,}000$, at most half a gap, about $3 \times 10^{-11}$~s. The error
is not zero, but about $5 \times 10^{-11}$~s, a billionth of the error of the chopped $23$-place tenth.

\textbf{19.} One tick of $0.1$~s chopped after $23$ binary places is short by $0.8 \times 2^{-23} \approx 9.54 \times
10^{-8}$~s; after $100$ hours, that is $3{,}600{,}000$ ticks, the clock is behind by $0.3433$~s, exactly the value
published in the report.
\end{solution}
